% Part: turing-machines
% Chapter: machines-computations
% Section: configuration

\documentclass[../../../include/open-logic-section]{subfiles}

\begin{document}

\olfileid{tur}{mac}{con}
\olsection{التهيئات والحسابات}

\begin{explain}
قد تتبعنا من قبل تهيئات آلة الزوجية. وما الشريط المتخيل والرأس والبرنامج إلا تصوير حدسي لحساب آلة تورنغ؛ أما حد الحساب الصوري فهو متتالية من \emph{التهيئات}. والتهيئة ثلاثية: سلسلة رموز تمثل محتوى الشريط في وقت من الحساب، وعدد يبين موضع الرأس، وحالة. فبهذه العناصر نضبط ما تحسبه الآلة~$\OLId{latin-looped}{M}$ على دخل معين.
\end{explain}

\begin{defn}[التهيئة]
\emph{تهيئة} آلة تورنغ $\OLId{latin-looped}{M}=\tuple{\OLId{latin-looped}{Q},\Sigma,\OLId{latin-ordinary}{q}_\OLMathNumeral{0},\delta}$ ثلاثية $\tuple{\OLId{latin-looped}{C},\OLId{latin-ordinary}{m},\OLId{latin-ordinary}{q}}$ تستوفي:
\begin{enumerate}
\item $\OLId{latin-looped}{C} \in \Sigma^*$ سلسلة منتهية من رموز $\Sigma$؛
\item $\OLId{latin-ordinary}{m} \in \Nat$ وأصغر من $\len{\OLId{latin-looped}{C}}$؛
\item $\OLId{latin-ordinary}{q} \in \OLId{latin-looped}{Q}$.
\end{enumerate}
ومعنى~$\OLId{latin-looped}{C}$ في التصوير محتوى الشريط، من أقصى خانة يسارًا إلى آخر خانة غير خالية أو بلغها الرأس من قبل. وأما $\OLId{latin-ordinary}{m}$ فرقم الخانة الممسوحة، إذ يبدأ الترقيم بـ$\OLMathNumeral{0}$ لأقصى اليسار؛ وأما $\OLId{latin-ordinary}{q}$ فالحالة الراهنة.
\end{defn}

\begin{explain}
دخل الآلة سلسلة رموز، ترمز عادة عددًا على وجه ما. والتهيئة الابتدائية هي وضع الآلة عند الشروع في حساب ذلك الدخل: في أول الشريط علامة طرفه، ويليها الدخل مباشرة إلى اليمين؛ والرأس يمسح أول خانة من الدخل، وهي التي تلي علامة الطرف الأيسر؛ والحالة هي الحالة الابتدائية المعينة~$\OLId{latin-ordinary}{q}_\OLMathNumeral{0}$.
\end{explain}

\begin{defn}[التهيئة الابتدائية]
إذا شغلت $\OLId{latin-looped}{M}$ على الدخل $\OLId{latin-looped}{I} \in \Sigma^*$، فـ\emph{تهيئتها الابتدائية}
\[
\tuple{\TMendtape \frown \OLId{latin-looped}{I} \frown \TMblank, \OLMathNumeral{1}, \OLId{latin-ordinary}{q}_\OLMathNumeral{0}}.
\]
\end{defn}

\begin{explain}
نستعمل~$\frown$ رمزًا لـ\emph{الوصل}. فسلسلة الدخل تتلو علامة الطرف من جهة اليمين بلا خانة فاصلة.
\end{explain}

\begin{defn}
التهيئة $\tuple{\OLId{latin-looped}{C},\OLId{latin-ordinary}{m},\OLId{latin-ordinary}{q}}$ \emph{تنتج $\tuple{\OLId{latin-looped}{C}',\OLId{latin-ordinary}{m}',\OLId{latin-ordinary}{q}'}$ في خطوة واحدة} بحسب~$\OLId{latin-looped}{M}$ إذا وفقط إذا اجتمعت الشروط:
\begin{enumerate}
\item رمز الموضع $\OLId{latin-ordinary}{m}$ من $\OLId{latin-looped}{C}$ هو $\sigma$؛
\item تعليمات $\OLId{latin-looped}{M}$ تعطي $\delta(\OLId{latin-ordinary}{q},\sigma)=
  \tuple{\OLId{latin-ordinary}{q}',\sigma',\OLId{latin-looped}{D}}$؛
\item رمز الموضع $\OLId{latin-ordinary}{m}$ من $\OLId{latin-looped}{C}'$ هو $\sigma'$؛
\item وتتحقق إحدى الحالات:
\begin{enumerate}
\item $\OLId{latin-looped}{D}=\OLId{latin-looped}{L}$، وعندئذ $\OLId{latin-ordinary}{m}'=\OLId{latin-ordinary}{m}-\OLMathNumeral{1}$ إن كان $\OLId{latin-ordinary}{m}>\OLMathNumeral{0}$، وإلا $\OLId{latin-ordinary}{m}'=\OLMathNumeral{0}$؛ أو
\item $\OLId{latin-looped}{D}=\OLId{latin-looped}{R}$ و$\OLId{latin-ordinary}{m}'=\OLId{latin-ordinary}{m}+\OLMathNumeral{1}$؛ أو
\item $\OLId{latin-looped}{D}=\OLId{latin-looped}{N}$ و$\OLId{latin-ordinary}{m}'=\OLId{latin-ordinary}{m}$؛
\end{enumerate}
\item إن كان $\OLId{latin-ordinary}{m}'=\len{\OLId{latin-looped}{C}}$، كان $\len{\OLId{latin-looped}{C}'}=\len{\OLId{latin-looped}{C}}+\OLMathNumeral{1}$ ورمز الموضع $\OLId{latin-ordinary}{m}'$ من $\OLId{latin-looped}{C}'$ هو~$\TMblank$؛ وإلا $\len{\OLId{latin-looped}{C}'}=\len{\OLId{latin-looped}{C}}$؛
\item لكل $\OLId{latin-ordinary}{i}$ يحقق $\OLId{latin-ordinary}{i}<\len{\OLId{latin-looped}{C}}$ و$\OLId{latin-ordinary}{i}\neq \OLId{latin-ordinary}{m}$ تبقى $\OLId{latin-looped}{C}'(\OLId{latin-ordinary}{i})=\OLId{latin-looped}{C}(\OLId{latin-ordinary}{i})$.
\end{enumerate}
\end{defn}

\begin{defn}\ollabel{defn:run-output}
\emph{تشغيل $\OLId{latin-looped}{M}$ على الدخل~$\OLId{latin-looped}{I}$} متتالية منتهية أو لا نهائية $\OLId{latin-looped}{C}_\OLId{latin-ordinary}{i}$ من تهيئات $\OLId{latin-looped}{M}$؛ أولها $\OLId{latin-looped}{C}_\OLMathNumeral{0}$ التهيئة الابتدائية لـ$\OLId{latin-looped}{M}$ على~$\OLId{latin-looped}{I}$، وكل $\OLId{latin-looped}{C}_\OLId{latin-ordinary}{i}$ تنتج $\OLId{latin-looped}{C}_{\OLId{latin-ordinary}{i}+\OLMathNumeral{1}}$ في خطوة واحدة متى وجدت $\OLId{latin-looped}{C}_{\OLId{latin-ordinary}{i}+\OLMathNumeral{1}}$. وإذا انتهت المتتالية، لم تنتج تهيئتها الأخيرة تهيئة أخرى في خطوة واحدة.

وتكون $\OLId{latin-looped}{M}$ \emph{متوقفة على الدخل $\OLId{latin-looped}{I}$ بعد $\OLId{latin-ordinary}{k}$ خطوات} إذا انتهى تشغيلها إلى $\OLId{latin-looped}{C}_\OLId{latin-ordinary}{k}=\tuple{\OLId{latin-looped}{C},\OLId{latin-ordinary}{m},\OLId{latin-ordinary}{q}}$، وكان رمز الموضع $\OLId{latin-ordinary}{m}$ من~$\OLId{latin-looped}{C}$ هو $\sigma$، ولم تكن $\delta(\OLId{latin-ordinary}{q},\sigma)$ معرفة. وحينئذ \emph{خرج}~$\OLId{latin-looped}{M}$ على~$\OLId{latin-looped}{I}$ هو~$\OLId{latin-looped}{O}$، حيث $\OLId{latin-looped}{O}$ سلسلة لا تنتهي بـ~$\TMblank$ وتحقق $\OLId{latin-looped}{C}=\TMendtape\concat \OLId{latin-looped}{O}\concat\TMblank^\OLId{latin-ordinary}{j}$ لبعض $\OLId{latin-ordinary}{j} \in \Nat$. ومعنى $\TMblank^\OLId{latin-ordinary}{j}$ سلسلة من $\OLId{latin-ordinary}{j}$ رموز كلها $\TMblank$.
\end{defn}

\begin{explain}
فخرج~$\OLId{latin-looped}{O}$ للآلة~$\OLId{latin-looped}{M}$ إما منته برمز غير~$\TMblank$، وإما أن يكون $\OLId{latin-looped}{O}$ خاليًا؛ وهذه الحالة الثانية إذا كان الشريط عند الزمن~$\OLId{latin-ordinary}{k}$ كله~$\TMblank$، سوى $\TMendtape$ في أقصى اليسار.
\end{explain}

\end{document}
